\subsection{Reading the Built-in Protocol Sketches}
\label{sec:builtin-reading}

Inheritance used an environment, scripts, and a queue. Before reading
another metaobject generator, inspect those components in
\filepath{examples/misc/built-ins.abclr2}. It describes environments,
script objects, and queues in familiar object notation. It is a
\emph{protocol sketch}, not a replacement runtime to load:
angle-bracket placeholders, spelling mistakes such as \code{Nmae}
and \code{Vlaue}, an extra bracket, and the \code{:bound} fallback
prevent that interpretation. Consult \filepath{src/primary-env.lsp}
and \filepath{src/primary-queue.lsp} for the supplied implementations.

An environment is an object representing variable bindings. Its local
frame is an association list: \code{((n . 2) (limit . 10))}
binds \code{n} to 2. A second environment, \code{lexEnv},
represents the enclosing lexical scope.

\wholelisting[firstline=9,lastline=20]{examples/misc/built-ins.abclr2}

Trace a lookup of \code{limit}: \code{assoc} finds its pair, and
\code{cdr} supplies 10 to the reply destination. If no local pair exists,
the request goes outward, preserving \code{@ R}.
A local binding whose value is \code{nil} still counts:
the test examines the pair, not its value. Assignment by \code{:set}
updates the nearest existing binding; it does not necessarily create a
new local binding. The actual implementation also accepts
\code{[:set-local Name Value]}, updating or creating a binding in
this frame. The evaluator uses this when naming a new object.
\code{:alist} and \code{:lexEnv} inspect the frame and its parent.
The global environment ends the chain and reports an unbound lookup.
The sketch's outer \code{:bound} request is a typo, not a second
protocol: the implementation forwards \code{:boundp}.

A script object separates two jobs usually hidden by a handler
declaration. \code{[:match Message Reply Sender]} returns either
\code{nil} for failure or \code{[Bindings ExecutableExpression]} for
success. The bindings are an association list; the executable expression
may be \code{[:script ScriptObject]}. The evaluator then asks that object
to execute with \code{[:exec Env Id Gid Eval] @ C}. Here \code{Env}
contains the bindings, \code{Id} is the current metaobject identity,
\code{Gid} the group manager, \code{Eval} the evaluator, and \code{C}
the execution continuation. Pattern and body inspection are separate
requests, not commands to run the script.

The queue sketch stores arrivals by \code{cons} at the front and removes
the last element. Arrivals A, then B produce \code{(B A)}, but dequeue
returns A first: this is \emph{first in, first out} (FIFO), not a stack. Its \code{waiters}
list similarly retains reply destinations for dequeue requests made
while empty. Such a request need not stop the queue from accepting a
later enqueue. The built-in queue uses lower-level machinery instead;
do not infer that its storage representation is this pair of lists.
It additionally provides \code{:listify}, used by the metaobjects below.

\paragraph{Check your understanding.}
An inner frame is \code{((n . nil))}; its parent binds \code{n} to 7.
What does looking up \code{n} return? \emph{Answer:} \code{nil},
because the local binding shadows the parent's binding.
Would storing A with \code{cons} make it the next queue item?
\emph{Answer:} only if no older item exists; removal is from the
opposite end.

\subsection{A Script Object with a Sender-sensitive Match}
\label{sec:user-script-reading}

The sketch separated script matching from execution; now follow an
executable script object in \filepath{examples/misc/user-script.abclr2}.
Loading it creates four objects; the installation commands at the end
are inside a block comment and must be issued separately.
\code{Monkey} initially replies \emph{Scratch!} to any message.
\code{Monkey-Showman} and \code{Poor-Audience} each handle \code{:try}
by sending \code{:play-trick} to Monkey and forwarding their caller's
reply destination. Monkey sees different \emph{senders}
even though the requests have the same content.

The fourth object, \code{monkey-script}, implements the script-object
contract. It is not Monkey's application handler until installed
in Monkey's script set. For \code{:play-trick}, its match method tests
\code{(eq Sender Monkey-Showman)}. On success, it returns two elements:

\begin{codeblock}
\begin{lstlisting}
[ [ ['R . Reply] ] [:script Me] ]
; First element: association list binding R to Reply.
; Second element: expression requesting this script's execution.
\end{lstlisting}
\end{codeblock}

These are nested lists, not three bindings. \code{R} is quoted
because it is a binding name; \code{Reply} is evaluated because its value
is the destination being bound. Here \code{Me} is the script object,
not Monkey or the Showman.
For a different sender the match method replies \code{nil}.
The next catch-all \code{:match} method also replies \code{nil};
without it, script search could wait rather than move on to
Monkey's original catch-all handler.

\wholelisting[firstline=51,lastline=66]{examples/misc/user-script.abclr2}

The execution method asks the environment for \code{R}. Its continuation
sends \emph{Hansei} to that application reply destination.
Separately it sends \code{nil} to \code{C}, completing the evaluator's
request. These destinations have different meanings.
There is an ordering boundary: the lookup uses \code{<=}, not
\code{<==}. The source does not wait for its continuation to deliver
\emph{Hansei} before sending completion to \code{C}.
Completion is not a receipt from the application caller.

Try these commands in a fresh session, preferably with a scratch copy:

\begin{codeblock}
\begin{lstlisting}
(inc "examples/misc/user-script" :use-source t)
[Monkey-Showman <== :try] ; "Scratch!"
[Poor-Audience <== :try]  ; "Scratch!"
[[meta Monkey] <= [:add-script monkey-script]]
[[meta Monkey] <== :scripts] ; inspect installation
[Monkey-Showman <== :try]
[Poor-Audience <== :try]
\end{lstlisting}
\end{codeblock}

The final replies are \emph{Hansei} and \emph{Scratch!}.
Installation puts the new script first. It handles the Showman's request;
its failed match for Audience lets search reach Monkey's original handler.
Removing a script with \code{[:remove-script message]} removes only the
first script matching that message with \code{nil} reply and sender fields,
not every matching script or a script selected by identity. The custom
sender test here matters: a removal request with a \code{nil} sender does
not select \code{monkey-script}. Script-list changes affect later dispatch,
not a handler already being evaluated.
\code{:add-script} does not send a reply: changing its arrow to
\code{<==} waits for a reply which the protocol never promises.

\paragraph{Check your understanding.}
How can the script distinguish Showman and Audience without adding a
field to \code{:play-trick}? \emph{Answer:} matching receives the sender
as well as content. Why must a failed match reply?
\emph{Answer:} search needs an explicit failure before trying another
candidate.

\subsection{MCP: Following a Meta-circular Interpreter}
\label{sec:mcp-reading}

We have followed individual environment and script protocols. Now combine
them: \filepath{examples/misc/MCP.abclr2} expresses a metaobject generator,
an evaluator, and a group manager in ABCL/R2 itself.
\emph{Meta-circular} means that the interpreter is described using
the language it interprets. It does not mean this file removes all
dependencies on host Lisp and existing kernel objects.
Read it in three passes: the group manager at the bottom, the generated
metaobject, and finally the evaluator's expression cases.

The ordinary object \code{Meta-circular} implements the group-manager
protocol. Its state selects \code{MC-meta-gen} and \code{MC-eval}.
Its \code{:new} handler delegates to the generator, adding the
evaluator and \code{Me} as group manager.
\code{:replace-meta-gen} and \code{:replace-evaluator} change these
configuration fields, not necessarily every previously created object.
\code{:new-group} delegates to the existing group machinery and sends
\code{:initialize} to the resulting manager.

The generator receives state specifications, lexical environment,
scripts, evaluator, and group manager. Extra arguments are warned about
and ignored. It creates a metaobject holding queue, state environment,
script set, evaluator, manager, and a mode initially \code{:dormant}.
Initializer expressions are evaluated by \code{:do-evlis};
\code{pairlis} pairs variable names with the returned values.
An environment extending \code{LexEnv} becomes the state, and the
generator replies with the metaobject.

For application traffic, follow this cycle:

\wholelisting[firstline=33,lastline=43]{examples/misc/MCP.abclr2}

\begin{enumerate}
\item \code{[:message Message Reply Sender]} enqueues the three fields.
If dormant, the metaobject becomes active and sends itself \code{:begin}.
\item \code{:begin} dequeues one request, finds its script, and extends
the state environment with match bindings.
\item The evaluator executes the body. Its completion continuation sends
\code{:end} to the metaobject.
\item \code{:end} starts another queued request or returns to dormant mode.
\end{enumerate}

Application replies use \code{Reply}; advancement of the cycle uses the
execution continuation. Reflective handlers expose state, scripts, queue, and
evaluator. The one-field \code{[:find-script Message]} query supplies
\code{nil} for reply and sender; ordinary dispatch uses the actual fields.

In \code{MC-eval}, \code{Exp} is a tagged representation, not necessarily
original source. \code{:compiled} invokes a compiled function with the
context. \code{:variable} looks up ordinary names in \code{Env};
\code{Me} and \code{Group} yield \code{[den Id]} and \code{Gid}.
\code{:self-evaluate} invokes a supplied function.
\code{:script} invokes script-object execution.
An unknown representation warns; this is not an unrestricted Common Lisp
evaluator.

The send cases expose the difference between reply and completion:

\begin{longtable}{L{0.25\textwidth}L{0.65\textwidth}}
\textbf{Expression tag} & \textbf{Completion and delivery}\\\hline
\endhead
\code{:send-past} & Sends \code{nil} to \code{C}, then forwards the
request to the target's metaobject with its separate \code{Reply} field.\\
\code{:send-future} & Sends \code{nil} to \code{C}, obtains the future's
reply destination, and forwards the request with that destination.\\
\code{:send-now} & Uses \code{C} as the application's reply destination;
the target's reply supplies the expression's result.\\
\end{longtable}

These describe the non-null-target paths. A now-send to \code{nil}
does not reply in this definition.
\code{:object-def} chooses the specified group or the current one,
creates an object, and, if named, installs its denotation through
\code{:set-local}. An explicit generator receives the group's evaluator;
otherwise creation goes through the group manager.
\code{:group-def} creates and optionally binds a group manager.

\code{:do-progn} evaluates expressions sequentially, returning the last
result, or \code{nil} for an empty body.
\code{:do-evlis} evaluates sequentially and reconstructs a list of
results. Both use continuations to wait for each expression.

Here is a small teaching probe of the interpreter, not another original example:

\begin{lstlisting}
(inc "examples/misc/MCP" :use-source t)
[object mc-counter
  (group Meta-circular)
  (state [n := 0])
  (script
    (=> :increment [n := (+ n 1)] !n)
    (=> :value !n))]
[mc-counter <== :increment] ; 1
[mc-counter <== :increment] ; 2
[mc-counter <== :value]     ; 2
[[meta mc-counter] <== :evaluator] ; MC-eval
\end{lstlisting}

All four observations pass on the current image. They establish this
creation/state/reply path, not coverage of every evaluator case.

\paragraph{Check your understanding.}
Why can a past-send complete before the receiver answers?
\emph{Answer:} its evaluator sends \code{nil} to \code{C} independently
of the application reply. A now-send instead lets the receiver answer
directly to \code{C}.

\subsection{Meta-meta: Reflecting on the Message Interpreter}
\label{sec:meta-meta-reading}

\filepath{examples/misc/meta-meta.abclr2} takes the queue/dispatch
cycle one level higher. Its application object, \code{twice},
replies with \code{(+ x x)}. The interesting change is which objects
interpret its requests.

Read the four numbered source steps in order.
\code{the-meta-meta-gen} is placed in
\code{(group-of (group-of group))}, and \code{the-meta-gen} in
\code{(group-of group)}. Each \code{group-of} selects the group of
its argument, climbing the group relationship in this setup;
\code{group} is the current-group pseudo-variable.
The second generator creates a metaobject explicitly configured with
\code{(meta-gen the-meta-meta-gen)}. That object's own execution
is therefore interpreted by a custom meta-metaobject.
The third step creates \code{a-group} using \code{the-meta-gen};
the fourth puts \code{twice} in that group.

Keep three identities distinct:

\begin{lstlisting}
twice                  ; application object
[meta twice]           ; interprets application messages
[meta [meta twice]]    ; interprets interpreter messages
\end{lstlisting}

For \code{[twice <== 3]}, the first interpreter receives a
\code{:message} envelope with application content 3.
The higher interpreter sees messages involved in executing that
interpreter, including its begin/end cycle.
The generators retain the familiar state initialization, matching,
environment extension, and completion protocol.
Compare their repeated bodies with \code{MC-meta-gen}: logging is in
the generated objects' \code{:message} handlers, not merely in the
factories' creation handlers.

The actual labels are \emph{Meta-gen receievs message} and
\emph{Meta-meta receievs message}; spelling is the original's.
Initialization also produces traffic. Not every log line is another
application request, and numeric object identifiers are not stable.

\begin{codeblock}
\begin{lstlisting}
(inc "examples/misc/meta-meta" :use-source t)
[twice <== 3] ; 6, with meta-level diagnostic output
\end{lstlisting}
\end{codeblock}

The commented historical transcript refers to \code{foo}, although
the file defines \code{twice}, and one logging label differs.
Treat it as an old illustration, not the literal session to reproduce.
The isolated check confirms 6, both actual logging markers, and
distinct identities at all three levels. It does not require a fixed
log order or number of lines.

\paragraph{Check your understanding.}
Does a meta-metaobject make \code{twice} a subclass of its interpreter?
\emph{Answer:} no. These are separate objects with an execution-control
relationship, not an inheritance hierarchy.

\paragraph{Reproducing the checks.}
Run \code{make -C tutorial reflection-check}. Eighteen checks use
distributed copies of \code{user-script}, \code{MCP}, and \code{meta-meta};
Their metaobjects retain distinct application and execution-group bindings.
Inputs, translations, and logs stay in \filepath{build/reflection-tests/}.
An additional 38 checks exercise addition, first-match removal, duplicate
scripts, reinstallation, retained state, and queued replies on both
primary dispatch paths. They guard script-list/cache consistency, not
arbitrary concurrent modification of an executing handler.
